\begin{longtable}{@{}l p{12.5cm}@{}}
  \caption{Skills per-instance answer comparison ($50\% \leq \Delta < 75\%$) for the 9 most divergent items: same question, each model's answer and reasoning.}\label{tab:cmp_skills_pos50} \\
  \toprule
  \endfirsthead
  \toprule
  \endhead
  \multicolumn{2}{@{}l}{\textbf{85580} \quad $\Delta = +70\%$} \\
  \textbf{Question} & You are an evaluator using a formative assessment system that analyzes student-submitted concept maps in real time to identify understanding of core concepts, connection errors, and missing supporting concepts in a Special Education / Career and Vocational Skills course. Given a hypothetical student concept map for a module titled "Workplace Readiness for Students with Autism Spectrum Disorder (ASD)" that includes the following nodes and links: "Social CuesWorkplace Communication"; "Task ChunkingIndependent Task Completion"; "Sensory RegulationOn-the-Job Behavior"; and an isolated node "Self-Advocacy" with no links, evaluate the map's conceptual strengths and weaknesses, justify which missing or incorrect connections most critically impair the student's readiness model, and design a prioritized plan of corrective instructional feedback and two concrete in-class activities the instructor should use immediately to remediate the highest-priority gaps and reinforce internalization of workplace values (reliability, responsibility, self-advocacy). Your plan should include (a) the single most important connection to add or fix and why, (b) three tailored feedback messages (brief, actionable) that the system can deliver in real time to the student, and (c) two activity descriptions (including objectives, brief procedures, and expected student behaviors/outcomes) that promote integration of concepts and value internalization. Be specific and justify your choices with respect to how they address conceptual misconceptions and support long-term characterizing of workplace values. \\
  \textbf{Reference} & Evaluation of the concept map's strengths and weaknesses: \newline Strengths: The student correctly identifies several relevant constructs for workplace readiness for students with ASD: social cues, task chunking, sensory regulation, workplace communication, independent task completion, and on-the-job behavior. The existing links show an emerging understanding that specific strategies (task chunking, sensory regulation) influence performance outcomes (independent task completion, on-the-job behavior), which is an appropriate causal structure. \newline Weaknesses: The map omits critical links connecting social-cognitive skills and metacognitive/self-determined skills to workplace outcomes. "Workplace Communication" is connected only from "Social Cues," but there is no explicit connection from "Self-Advocacy" to any outcome or to communication. "Self-Advocacy" being isolated is the most important omission because self-advocacy directly mediates accommodations, help-seeking, and integrity of responsibilities---key components of reliability and responsibility. Additionally, there is no explicit link showing that emotional or sensory regulation supports social interactions (i.e., sensory regulation workplace communication) or that self-advocacy interacts with task chunking (e.g., asking for task breakdowns). The map also lacks a node representing "Accommodation/Support Requests" or "Workplace Reliability" as an outcome that synthesizes these inputs, which weakens the student's model of how values translate into measurable behaviors. \newline Most critical missing/incorrect connection (single most important fix) and justification: \newline Add and prioritize the connection: "Self-AdvocacyAccommodation/Support RequestsIndependent Task Completion" and directly link "Self-AdvocacyWorkplace Communication." Justification: Self-advocacy is the mechanism by which a worker with ASD requests needed supports (clarification, schedule adjustments, sensory breaks), which then enables reliable task completion and demonstrates responsibility. Without this connection, the map treats self-advocacy as optional rather than central to workplace success and value internalization. Fixing this link explicitly ties an internalized value (making responsible requests when needed) to observable outcomes (on-time task completion, appropriate communication), which supports both practical readiness and characterizing workplace values. \newline Three tailored, real-time feedback messages the system should deliver (brief and actionable): \newline \textit{[\ldots\ truncated \ldots]} \newline - The feedback messages are immediate and actionable, enabling students to revise their conceptual models during the learning episode. The activities combine practice, social modeling, and reflection---essential components for characterizing values---so students repeatedly choose responsible behaviors until they become part of their professional identity. \newline In sum, by adding and emphasizing the connection from self-advocacy to accommodation/support requests and to workplace communication, and by using targeted real-time feedback plus role-play and scaffolded problem-solving activities, the evaluator supports both conceptual accuracy and the internalization of workplace values (reliability, responsibility, self-advocacy) in a measurable, practice-based way. \\
  \midrule
  \textbf{Kimi-K2.6} & score 2/10 \\
  & [1-2] The student's answer is severely incomplete, failing to address weaknesses, connection errors, missing concepts, or the required instructional plan. It demonstrates insufficient basics and a serious misunderstanding of the task requirements, fitting the lowest band. \\
  \textbf{Innospark-1T} & score 9/10 \\
  & Let me think through this carefully. \newline The question/task: You are an evaluator using a formative assessment system that analyzes student-submitted concept maps in real time to identify understanding of core concepts, connection errors, and missing supporting concepts in a Special Education / Career and Vocational Skills course. Given a hypothetical student concept map for a module ti \newline My reasoning: Formative Assessment Evaluation: Concept Map Analysis \newline Step-by-Step Reasoning \newline Step 1: Analyze the Current Concept Map Structure \newline Present Nodes and Links: \newline \textit{[\ldots\ truncated \ldots]} \newline - Activity 2: "Workplace Scenario Simulation" - Role-play to practice skill integration and value demonstration \newline These interventions address the root misconception (isolated skills) while building toward long-term internalization of workplace values through active, meaningful engagement with interconnected concepts. \\
  \specialrule{1.2pt}{8pt}{8pt}
  \multicolumn{2}{@{}l}{\textbf{39040} \quad $\Delta = +70\%$} \\
  \textbf{Question} & You are designing a multimodal instructional module for an advanced linguistics course that teaches students how to produce persuasive short documentary videos (3--5 minutes) that communicate sociolinguistic research findings to non-specialist audiences. Describe and justify your complete design choices across image, audio, and video modalities. Specifically, (1) identify the target audience and learning objectives; (2) select and justify the multimodal affordances (visual framing, shot types, graphic overlays, color palette, voice style, pacing, ambient sound, music, captions, and image/sound editing techniques) you will use to make complex linguistic concepts accessible without oversimplification; (3) explain how you will structure the 3--5 minute narrative (sequence of scenes, evidence presentation, and persuasive appeals) and tie each segment to specific multimodal decisions; (4) discuss how you will ensure ethical representation of communities and accurate portrayal of data (including consent, anonymization, and visualization choices); and (5) propose a brief evaluation plan to judge whether the video succeeded in communicating the research to the target audience (specify which qualitative and/or quantitative measures you would collect). Justify each decision with reference to multimodal communication principles and principles from linguistics (e.g., indexicality, register, speech accommodation). \\
  \textbf{Reference} & Target audience and learning objectives: The target audience is educated non-specialists aged 18--35 (undergraduates, policy interns, community organizers) with interest in language and society but limited technical training in linguistics. Learning objectives are: (1) convey the central research question and main empirical findings of a sociolinguistic study in plain language; (2) demonstrate how linguistic evidence supports the claims using clear multimodal examples; (3) foster critical reflection on social implications and ethical considerations. These objectives prioritize comprehension, evidentiary reasoning, and ethical awareness---appropriate for persuasive documentary style aimed at motivating informed action or reflection. \newline Multimodal affordances and justifications: \newline - Visual framing and shot types: Use a mix of medium close-ups for speaker credibility and intimacy (when introducing researchers or community members), wide establishing shots to situate community context, and extreme close-ups for visualizing micro-evidence (e.g., screenshots of transcribed utterances, spectrogram excerpts). Close-ups support ethos and pathos; wide shots provide contextual indexicality showing environment-linked language variation. Macro-to-micro alternation helps viewers move from general claims to concrete evidence. \newline - Graphic overlays and data visualization: Design simple, annotated visuals---color-coded utterance excerpts, timelines showing change over time, and minimally-annotated spectrograms with salient features highlighted. Use consistent visual metaphors (e.g., speech balloons linked to speakers) to map audio examples to text. Avoid dense charts; prefer annotated exemplars that preserve linguistic nuance. This balances accessibility with fidelity to data and leverages dual coding (visual + verbal) to scaffold comprehension. \newline - Color palette and typography: Use a restrained, high-contrast palette (two primary colors plus neutrals) to guide attention without distracting. Typography should use a legible sans-serif for captions and a monospace-like font for transcriptions to signal analytic content. Color and type cues create visual hierarchy and signal transitions between narrative, data, and interpretation. \newline - Voice style and onscreen narration: Adopt a conversational yet precise narratorial voice---moderate speaking rate, warm timbre, and clear enunciation---to maximize comprehensibility and trustworthiness. When presenting technical terms, use brief lay definitions and voice modulation to mark importance. This reflects accommodation: matching register to audience while retaining disciplinary accuracy. \newline \textit{[\ldots\ truncated \ldots]} \newline - Ethical linguistics: Transparency about transcription choices, anonymization, and community review respects epistemic responsibility and reduces the risk of extractive representation. \newline Together, these design choices ensure the short documentary is persuasive yet rigorous: it makes principled trade-offs between simplicity and fidelity, uses multimodal affordances to scaffold understanding, and embeds ethical safeguards to respect and accurately represent linguistic communities. \\
  \midrule
  \textbf{Kimi-K2.6} & score 2/10 \\
  & [1-2] The answer is severely incomplete, cutting off mid-sentence in section 2 and entirely missing sections 3-5. This results in limited learning effectiveness, no evidence of accessibility/ethical considerations, and failure to integrate the required educational design elements. \\
  \textbf{Innospark-1T} & score 9/10 \\
  & This is a comprehensive instructional design question about creating a multimodal educational module for advanced linguistics students. Let me work through each component systematically, drawing on principles from multimodal communication, documentary filmmaking, and sociolinguistics. \newline Let me structure my response to address all five required elements: \newline 1. Target audience and learning objectives \newline 2. Multimodal affordances with justification \newline 3. Narrative structure tied to multimodal decisions \newline 4. Ethical considerations \newline \textit{[\ldots\ truncated \ldots]} \newline - Media literacy and ethics \newline Let me develop each section thoughtfully. \\
  \specialrule{1.2pt}{8pt}{8pt}
  \multicolumn{2}{@{}l}{\textbf{113823} \quad $\Delta = +60\%$} \\
  \textbf{Question} & You are a research assistant tasked with producing a prioritized list of novel K-12 research ideas in Technology and Home Economics (ITEEA/STEL; FCS standards). Based on an automated literature scan that summarizes recent papers and highlights gaps, analyze three candidate research topics that emerged from the scan. For each topic: (1) briefly summarize the gap identified in the literature, (2) explain the potential educational and ethical values at stake (e.g., equity, student autonomy, cultural relevance, safety), (3) compare and prioritize the three topics in order of recommended pursuit, and (4) justify your prioritization by resolving value conflicts and describing how the top-priority study would balance practical feasibility, impact on students' affective outcomes (e.g., attitudes, motivation, value internalization), and ethical concerns. Provide a detailed, organized response that demonstrates clear reasoning for the chosen ordering and concrete steps the top-priority project should take to uphold the prioritized values. \\
  \textbf{Reference} & Topic summaries and prioritization \newline Topic A: Integrating maker-centered engineering projects into middle school family and consumer science (FCS) to promote student identity as creators \newline 1) Gap summary: Recent automated scans show many studies on maker education in general STEM but few that specifically integrate maker-centered engineering within FCS curricula. Existing work often treats technology and home economics separately; little research examines how maker activities in FCS influence students' sense of competence, identity as creators, and sustained interest in technical careers, particularly among girls and underrepresented groups. \newline 2) Values at stake: Equity (access for low-income schools to maker tools), student autonomy (choice in project topics), authenticity (cultural relevance of projects tied to home/family contexts), and safety (tools and materials). Tension may arise between providing advanced tools (promotes competence) and ensuring equitable distribution. \newline Topic B: Culturally responsive digital nutrition education using smartphone apps for diverse K-12 communities \newline 1) Gap summary: Literature shows many digital nutrition interventions but sparse research tailoring content culturally for diverse K-12 populations and measuring affective outcomes like value internalization toward healthy eating. There is also limited research on student and family engagement across cultural groups. \newline \textit{[\ldots\ truncated \ldots]} \newline Concluding rationale \newline Topic A offers a strong balance: it fills a tangible literature gap, aligns with curricular standards, promises measurable affective benefits (identity, motivation), and allows clear, ethically sensitive implementation strategies that mitigate equity and safety concerns. Prioritizing it enables scalable research with actionable outcomes for practitioners while maintaining respect for student values and community contexts. \\
  \midrule
  \textbf{Kimi-K2.6} & score 2/10 \\
  & The user wants a prioritized list of three novel K-12 research ideas in Technology and Home Economics (aligned with ITEEA/STEL and FCS standards). I need to: \newline 1. Analyze three candidate research topics that emerged from an (hypothetical) automated literature scan. \newline 2. For each topic: \newline - Summarize the gap identified in the literature. \newline - Explain potential educational and ethical values at stake (equity, student autonomy, cultural relevance, safety). \newline 3. Compare and prioritize the three topics in order of recommended pursuit. \newline \textit{[\ldots\ truncated \ldots]} \newline One final check: Are these truly "novel"? Yes, the specific integration points are under-researched. The repair topic is particularly timely with right-to-repair legislation. The cultural food tech topic addresses a real gap in how engineering design treats non-dominant knowledge. The neurodivergent AI topic is forward-looking but ethically fraught, making it a good third choice. \newline I will write the response now. \\
  \textbf{Innospark-1T} & score 8/10 \\
  & Let me think through this carefully. \newline The question/task: You are a research assistant tasked with producing a prioritized list of novel K-12 research ideas in Technology and Home Economics (ITEEA/STEL; FCS standards). Based on an automated literature scan that summarizes recent papers and highlights gaps, analyze three candidate research topics that emerged from the scan. For each topic: (1) brief \newline My reasoning: Research Prioritization Analysis: K-12 Technology and Home Economics \newline Methodology \newline This analysis synthesizes gaps identified in recent ITEEA/STEL and Family and Consumer Sciences literature, evaluating three emerging research topics through a multi-criteria framework. Topic Analysis \newline Topic 1: AI-Assisted Personalized Learning in Technology Education \newline \textit{[\ldots\ truncated \ldots]} \newline Conclusion \newline The STEM-Home Economics integration study represents the optimal balance of practical feasibility, positive affective impact, and manageable ethical considerations. By prioritizing this research, the field can build foundational knowledge about cross-disciplinary approaches before tackling more complex technological interventions. \\
  \specialrule{1.2pt}{8pt}{8pt}
  \multicolumn{2}{@{}l}{\textbf{115990} \quad $\Delta = +50\%$} \\
  \textbf{Question} & You are a research assistant given a list of recent peer-reviewed K--12 mathematics education papers focused on implementing Common Core State Standards for Mathematics (CCSS-M) in heterogeneous middle school classrooms. Using only the paper list (titles, authors, abstracts, and publication years), perform a focused literature review to (1) identify and categorize the principal research themes and methodologies represented in the corpus, (2) analyze gaps or under-researched aspects across those themes specifically related to equitable access and learning outcomes, and (3) propose two coherent, novel research questions that address the most critical gaps you identified. In your answer, (a) clearly describe how you grouped the papers into themes and why, (b) summarize the dominant methodological approaches and note any methodological limitations you observed, (c) provide evidence-based reasoning for each identified gap with references to examples from the paper list (cite paper titles), and (d) present two specific research questions (with brief rationale and suggested study design for each). Assume the paper list contains at least 12 items with varied focuses (curriculum alignment, teacher professional development, formative assessment, technology integration, differentiated instruction, student discourse, etc.). \\
  \textbf{Reference} & Grouping and thematic categorization \newline I began by reading titles and abstracts to identify recurring focal points and clustered the corpus into four principal themes based on topic similarity and research aims: (A) Curriculum Alignment and Instructional Materials (papers that study alignment of classroom tasks/materials to CCSS-M standards and fidelity of implementation), (B) Teacher Professional Development and Instructional Practices (studies on PD models, teacher beliefs, and enacted pedagogy), (C) Formative Assessment, Feedback, and Student Discourse (research on assessment-for-learning, questioning strategies, and classroom talk), and (D) Differentiation, Technology, and Equity (work addressing adaptive tools, differentiated instruction strategies, access issues, and underrepresented students). I grouped papers when abstracts emphasized the same primary problem (e.g., "alignment to standards" vs. "teacher learning") and when methods and outcome measures overlapped (e.g., classroom observations and lesson artifact analyses for implementation studies). This clustering highlights how research attention is distributed across policy-to-practice, teacher learning, assessment practices, and equity/technology concerns. \newline Dominant methodological approaches and limitations \newline Across themes, three methodological patterns dominate: (1) qualitative case studies and classroom observations (common in Teacher PD and Student Discourse papers), (2) mixed-methods designs combining surveys, pre/post assessments, and classroom artifacts (frequent in Curriculum Alignment and PD studies), and (3) quasi-experimental or small randomized controlled trials (RCTs) for technology and some formative assessment interventions. Methodological limitations I observed include: small and localized samples limiting generalizability (e.g., several PD studies sample 2--4 schools), short intervention durations that preclude claims about sustained effects (common in technology pilots), reliance on teacher self-report data without triangulation (some studies on alignment), and limited attention to heterogeneous student subgroups (many studies report average effects but do not disaggregate by prior achievement, English learner status, or socioeconomic status). Additionally, few studies combine fine-grained classroom discourse analysis with robust student learning measures, creating a gap between observed interactional mechanisms and measured outcomes. \newline Analysis of gaps related to equitable access and learning outcomes \newline From the thematic map and methods, three interrelated gaps emerge that affect equitable access and learning outcomes: \newline \textit{[\ldots\ truncated \ldots]} \newline Concluding synthesis \newline By categorizing the corpus into themes, reviewing dominant methods, and pinpointing recurring methodological blind spots, I identified actionable gaps---especially around differential effects for marginalized students and linking discourse mechanisms to outcomes. The two proposed research questions directly address these gaps with mixed-methods or experimental designs that emphasize subgroup analyses and scalability, aligning research priorities with equitable implementation of CCSS-M in heterogeneous middle school settings. \\
  \midrule
  \textbf{Kimi-K2.6} & score 4/10 \\
  & The user wants me to act as a research assistant and perform a focused literature review based on a hypothetical list of peer-reviewed K-12 mathematics education papers. However, the user hasn't actually provided the list of papers. They said "Assume the paper list contains at least 12 items..." but didn't include the list. \newline Wait, looking carefully at the prompt: "Using only the paper list (titles, authors, abstracts, and publication years)... Assume the paper list contains at least 12 items..." \newline Since no actual list was provided, I need to either: \newline 1. Ask for the list, or \newline 2. Create a plausible synthetic corpus based on the description and perform the review on that synthetic corpus. \newline Given the instructions say "Using only the paper list" and "Assume the paper list contains...", it seems the user wants me to work with a hypothetical/assumed list. The most helpful approach is to construct a representative synthetic corpus of 12+ papers that fits the description, explicitly state that I'm using this constructed representative corpus based on typical publications in this field, and then perform the analysis. \newline \textit{[\ldots\ truncated \ldots]} \newline Final check for conciseness. The user wants concise. I'll use bullet points and tight paragraphs. \newline Let's go. \\
  \textbf{Innospark-1T} & score 9/10 \\
  & Let me think through this carefully. \newline The question/task: You are a research assistant given a list of recent peer-reviewed K--12 mathematics education papers focused on implementing Common Core State Standards for Mathematics (CCSS-M) in heterogeneous middle school classrooms. Using only the paper list (titles, authors, abstracts, and publication years), perform a focused literature review to (1) ide \newline My reasoning: Focused Literature Review: CCSS-M Implementation in Heterogeneous Middle School Classrooms \newline Methodology for Thematic Analysis \newline I employed a systematic qualitative content analysis approach to categorize the papers. Grouping Process: \newline - Initial open coding of key concepts in each paper's title and abstract \newline \textit{[\ldots\ truncated \ldots]} \newline Summary \newline This literature review reveals that while the field has made progress in understanding CCSS-M implementation mechanics, significant gaps remain in understanding equity implications. The proposed research questions aim to advance the field by centering equity as a primary analytical lens and employing methodological approaches that can capture complexity across multiple levels of the educational system. \\
  \specialrule{1.2pt}{8pt}{8pt}
  \multicolumn{2}{@{}l}{\textbf{3403} \quad $\Delta = +50\%$} \\
  \textbf{Question} & You are an instructional designer conducting a task analysis for a middle school music lesson whose instructional goal is: 'Students will be able to perform a three-measure melodic phrase in 4/4 meter using correct rhythm, articulation, and expression.' Using Gagn's learning taxonomy, classify this instructional goal into the appropriate learning domain and subtypes (e.g., intellectual skills, verbal information, cognitive strategies, motor skills, attitudes). Then identify and justify the main steps (learning events or component tasks) you would include in a sequenced task analysis to teach this goal. For each main step, specify the prerequisite knowledge or skill, the observable performance outcome, and one concrete instructional activity or practice task that applies the prerequisite toward the outcome. Provide a clear rationale (brief) linking each step to Gagn's taxonomy and to the cognitive Applying level (use of knowledge in a new situation). \\
  \textbf{Reference} & Classification and overall rationale: \newline The instructional goal belongs primarily to the motor skills domain with a strong component of Gagn's "intellectual skills" and some elements of verbal information and cognitive strategies. Performing a melodic phrase requires fine motor coordination (motor skills), applying musical rules (intellectual skills---discriminations and rule application), remembering notation and expressive markings (verbal information), and using rehearsal/monitoring strategies (cognitive strategies). Because the emphasis is on applying known musical conventions to produce a specific performance, this aligns with the Applying level of Bloom's taxonomy: students use existing knowledge (rhythmic values, articulation symbols, expressive directions) in a new performance context. \newline Main steps (sequenced task analysis) with prerequisites, observable outcomes, activities, and rationale: \newline 1) Step: Decode notation and expressive markings for the three-measure phrase \newline - Prerequisite: Ability to read basic rhythmic values (quarter, eighth, half, whole), pitch names on the instrument/voice range, and common articulation/expression symbols (staccato, legato, accent, dynamics symbols). \newline - Observable performance outcome: Student accurately identifies and verbally labels rhythms, pitches, articulations, and dynamics in the printed phrase without performing it (e.g., "Measure 1: quarter, two eighths; pitch sequence C-D-E; staccato on beat 2; mezzo-forte"). \newline \textit{[\ldots\ truncated \ldots]} \newline Integration to Applying level and summary: \newline Each step is sequenced from simpler cognitive/verbal tasks (decoding notation) to isolated skill practice (rhythm, pitch), then to integrated motor performance, articulation/expression refinement, and finally transfer to realistic performance---this sequencing follows Gagn's emphasis on prerequisite events and task complexity. The design targets Applying by requiring learners to use prior knowledge (notation, rhythm conventions, instrument technique) in progressively more complex, realistic performance situations, demonstrating their ability to transfer and implement knowledge rather than merely recall it. \\
  \midrule
  \textbf{Kimi-K2.6} & score 4/10 \\
  & The user wants a task analysis for a middle school music lesson based on Gagn's learning taxonomy. \newline First, I need to classify the instructional goal: "Students will be able to perform a three-measure melodic phrase in 4/4 meter using correct rhythm, articulation, and expression." \newline Gagn's taxonomy includes: \newline 1. Verbal Information \newline 2. Intellectual Skills (discriminations, concrete concepts, defined concepts, rules, higher-order rules/problem solving) \newline 3. Cognitive Strategies \newline \textit{[\ldots\ truncated \ldots]} \newline I'll make sure the rationale is brief but covers both required links. \newline Okay, I'm ready to write. \\
  \textbf{Innospark-1T} & score 9/10 \\
  & Let me analyze this instructional design problem systematically. \newline **Step 1: Analyze the Goal Using Gagn's Taxonomy** \newline The goal is: "Students will be able to perform a three-measure melodic phrase in 4/4 meter using correct rhythm, articulation, and expression." \newline Breaking this down: \newline - "Perform" = physical execution on an instrument or voice \newline - "Three-measure melodic phrase" = specific musical content \newline \textit{[\ldots\ truncated \ldots]} \newline **Step 4: Develop Detailed Task Analysis** \newline Let me create a comprehensive task analysis with prerequisites, outcomes, and activities for each step. \\
  \specialrule{1.2pt}{8pt}{8pt}
  \multicolumn{2}{@{}l}{\textbf{49321} \quad $\Delta = +50\%$} \\
  \textbf{Question} & You are a teaching assistant preparing an individualized summary and learning path for a middle school music student based on their recent performance and reflection data. Analyze the student's profile below and then: (1) identify and explain three affective learning goals (valuing level) that should be prioritized for this student, (2) propose two specific classroom activities that would help the student develop each goal (six activities total), and (3) describe how you would monitor and document the student's growth in each affective goal over a six-week period. \newline Student profile: The student participates in choir and beginning piano. Their practice logs show consistent technical practice on scales but limited time on song repertoire. In class, they volunteer occasionally but rarely choose solos; they rate their confidence performing for peers as 2/5 and say they "enjoy music but worry about making mistakes." Their reflections show that they value teamwork but are unsure how their individual contributions affect the group sound. Teacher observations note steady technical improvement, hesitancy to take risks, and positive peer relationships. \newline Provide a detailed, reasoned response that connects the proposed goals and activities to the student's affective needs and explains recordkeeping methods that would inform a personalized learning path for the instructor. \\
  \textbf{Reference} & Model Answer: \newline 1) Prioritized affective learning goals (valuing level) and rationale \newline - Goal A: Develop a stronger sense of personal musical worth and performance confidence. Rationale: The student reports low confidence (2/5) and hesitancy to take solos despite technical competence. Strengthening their valuing of their own musical contribution will increase willingness to perform and take risks. \newline - Goal B: Cultivate an understanding of their individual role within ensemble music and the value of individual contributions to group sound. Rationale: The student values teamwork but is unsure how their part affects the ensemble. Helping them see the worth of their role will deepen their motivation to engage in repertoire work beyond technical practice. \newline - Goal C: Foster a positive, risk-tolerant attitude toward rehearsal and public performance (valuing effort and learning from mistakes). Rationale: The student enjoys music but fears mistakes, which limits expressive choices and repertoire practice. Encouraging a growth-oriented attitude will help them reframe mistakes as learning opportunities and increase participation. \newline 2) Two specific classroom activities for each goal (six activities total), with explanation of how each activity targets the goal \newline \textit{[\ldots\ truncated \ldots]} \newline - Adaptive actions based on evidence: If confidence scores and voluntary performances rise, the instructor is advised to offer increasingly public performance opportunities. If the student still hesitates, adjust activities toward smaller, scaffolded exposure and increase positive peer feedback. If the student shows improved understanding of ensemble role, the path shifts to more complex part-leadership and arranging tasks to deepen musical ownership. \newline Overall, these goals, activities, and monitoring steps directly address the student's affective needs---building personal musical worth, clarifying ensemble contribution, and fostering risk-taking---while producing concrete data the instructor can use to tailor instruction and move the student along a personalized learning path. \\
  \midrule
  \textbf{Kimi-K2.6} & score 4/10 \\
  & The user wants a detailed, reasoned response for a middle school music student. I need to: \newline 1. Identify and explain three affective learning goals at the **valuing level** (Krathwohl's taxonomy: receiving, responding, valuing, organization, characterization). Valuing means the student sees worth in something and behaves consistently with that value. It ranges from acceptance of value to preference for a value to commitment. \newline 2. Propose two specific classroom activities per goal (six total). \newline 3. Describe how to monitor and document growth over six weeks. \newline Student profile analysis: \newline - Choir + beginning piano \newline \textit{[\ldots\ truncated \ldots]} \newline Also mention the technical improvement --- we can build on that strength. \newline Alright, ready to write. \\
  \textbf{Innospark-1T} & score 9/10 \\
  & Let me analyze this student's profile systematically to identify affective learning needs and develop appropriate goals and activities. \newline **Student Profile Analysis:** \newline 1. **Technical strengths**: Consistent practice on scales, steady technical improvement \newline 2. **Performance anxiety**: Rates confidence at 2/5, worries about mistakes, avoids solos \newline 3. **Social/team orientation**: Values teamwork, has positive peer relationships \newline 4. **Self-perception gap**: Unsure how individual contributions affect group sound \newline \textit{[\ldots\ truncated \ldots]} \newline - Portfolio documentation of growth \newline - Pre/post self-assessment measures \\
  \specialrule{1.2pt}{8pt}{8pt}
  \multicolumn{2}{@{}l}{\textbf{54243} \quad $\Delta = +50\%$} \\
  \textbf{Question} & You are tutoring an undergraduate chemistry student who understands basic acid--base concepts and is beginning to study buffer solutions. Provide a detailed explanation and worked example that applies the Henderson--Hasselbalch equation to design a buffer. Specifically: (a) Explain how the Henderson--Hasselbalch equation is derived from the acid dissociation equilibrium and under what assumptions it is valid. (b) Using that equation, design a 1.00 L buffer solution at pH 4.50 with total (acid + conjugate base) concentration of 0.100 M using acetic acid (pKa = 4.76). Calculate the required amounts (in moles and grams) of acetic acid (CH3COOH, molar mass = 60.05 g/mol) and sodium acetate trihydrate (CH3COONa3H2O, molar mass = 136.08 g/mol) to prepare the buffer. Show all steps, approximations, and explain how you check whether the Henderson--Hasselbalch assumptions are reasonable for this case. \\
  \textbf{Reference} & (a) Derivation and assumptions: \newline Start from the acid dissociation equilibrium for a weak acid HA: HA H+ + A. The acid dissociation constant Ka is defined as Ka = [H+][A]/[HA]. Taking negative base-10 logarithms and rearranging gives pKa = pH log([A]/[HA]) which can be rearranged to the Henderson--Hasselbalch equation: \newline pH = pKa + log([A]/[HA]). \newline This form follows directly from the definition of Ka: log([A]/[HA]) = pH pKa. \newline Key assumptions for using Henderson--Hasselbalch: \newline - Activities are approximated by concentrations (ideal dilute solution). This is acceptable at relatively low ionic strengths. \newline \textit{[\ldots\ truncated \ldots]} \newline - CH3COONa3H2O: 0.03548 mol 4.83 g \newline Dissolve in water and dilute to 1.00 L. Check final pH and adjust slightly if needed with small amounts of HCl or NaOH. \\
  \midrule
  \textbf{Kimi-K2.6} & score 5/10 \\
  & [5-6] Part (a) is accurate and well-explained, but the answer is abruptly cut off in part (b), missing the final concentrations, mole and gram calculations, and the crucial assumption checks. This significant omission limits understanding and leaves the worked example incomplete. \\
  \textbf{Innospark-1T} & score 10/10 \\
  & Let me work through this problem step by step. \newline First, I need to understand what's being asked: You are tutoring an undergraduate chemistry student who understands basic acid--base concepts and is beginning to study buffer solutions. Provide a detailed explanation and worked example that applies the Henderson--Hasselbalch equation to design a buffer. Specifically: (a) Explain how the Henderson--Hasselbalch equation is derived from the acid dissociation equilibri \newline Key values given: ,, 4.50, 0.100, 4.76, 1.00, 4.50 \newline My approach: \newline 1. Set Up the Henderson--Hasselbalch Equation \newline \$\$\textbackslash{}text\{pH\} = \textbackslash{}text\{p\}K\_a + \textbackslash{}log\textbackslash{}frac\{[\textbackslash{}text\{A\}\textasciicircum{}-]\}\{[\textbackslash{}text\{HA\}]\}\$\$ \newline \textit{[\ldots\ truncated \ldots]} \newline | Component | Moles | Mass | \newline This gives us the final answer: | \\
  \specialrule{1.2pt}{8pt}{8pt}
  \multicolumn{2}{@{}l}{\textbf{14308} \quad $\Delta = +50\%$} \\
  \textbf{Question} & You are preparing an environmental analysis report for a kindergarten English Language Arts program. Based on your analysis of the user's educational values, compatibility with existing educational systems, and the physical and social characteristics of the facility, identify and prioritize three core educational values that should guide program decisions. For each value: (a) explain why it is prioritized given the user's context and compatibility with existing systems; (b) describe one likely conflict that could arise between this value and another stakeholder value or system constraint; and (c) propose a practical strategy to resolve or mitigate that conflict in the classroom and facility design. Support your reasoning with specific examples tied to kindergarten ELA practices and the facility's features. \\
  \textbf{Reference} & After reviewing the user's stated educational values and the characteristics of the existing education system and facility, I prioritize these three core educational values for a kindergarten ELA program: child-centered playfulness, language-rich inclusivity, and safe routine consistency. Below I justify each choice, identify likely conflicts, and propose practical resolution strategies grounded in kindergarten ELA practices and the facility's features. 1) Child-centered playfulness --- Priority justification: The user values early childhood engagement and the existing system allows some curricular flexibility; play-based approaches align with developmental ELA outcomes (oral language development, emergent literacy) and are well suited to a kindergarten schedule. The facility's features---if there are open play areas, low shelving, and natural light---support hands-on centers, storytelling corners, and interactive language games. Conflict: A likely conflict is pressure from standardized expectations or rigid curriculum pacing in the broader education system that emphasizes measurable skills over free play. Teachers or administrators may feel required to allocate more time to formal literacy drills, reducing play opportunities. Resolution strategy: Integrate playful activities that intentionally target ELA objectives so play counts as assessment time. For example, convert a dramatic play corner into a "post office" where children write and dictate simple notes (letter formation, vocabulary) and teachers use anecdotal records to document language use. Schedule short, focused explicit instruction blocks (10--15 minutes) for phonological awareness, paired with extended playful centers where those same skills are practiced in authentic contexts. Use the facility's open area for rotating play-based literacy stations to demonstrate alignment with system standards while preserving play-centered pedagogy. 2) Language-rich inclusivity --- Priority justification: The user emphasizes respecting children's home languages and diverse cultural backgrounds; the existing education system may require English outcomes but can accommodate multilingual supports. The facility's social characteristics (multicultural community, family involvement) make inclusivity vital for engagement and equitable language development. Conflict: Tension may arise between inclusive multilingual practices and stakeholders who prioritize rapid English-only proficiency or who lack resources for bilingual materials. Parents or administrators seeking quick measurable English gains may view time spent on home-language use as delaying progress. Resolution strategy: Adopt additive bilingual approaches that validate home languages while scaffolding English. Practically, create dual-language book corners and picture labels in children's home languages plus English, and involve families in creating word walls from home-language story contributions. Use echo-reading and shared book routines where teachers model English while encouraging children to retell stories in their home language, then bridge to English vocabulary. Leverage facility spaces like a family-resource corner for bilingual materials and parent-led story sessions, demonstrating how multilingual practices support, rather than hinder, English acquisition through transfer and increased language confidence. 3) Safe routine consistency --- Priority justification: Young learners thrive on predictable routines for emotional security, which supports affective engagement and readiness for literacy learning. If the facility has limited transitional spaces or a large class size, consistent routines reduce chaos and maximize instructional time. The existing system may require attendance, reporting, and behavior expectations that are better met when children have stable daily structures. Conflict: Strict routines can conflict with the child-centered value of flexibility and exploratory play; staff turnover or scheduling constraints may also make maintaining consistent routines challenging. Resolution strategy: Co-design predictable yet flexible routines with staff and families. For example, establish a consistent morning literacy ritual (welcome song, shared reading, choice time) that always includes a language-rich component but allows children to choose centers that reflect evolving interests. Use visual schedules and classroom signals placed at child eye level to support transitions, and configure facility zones (reading nook, writing table, sensory area) so that even when teacher schedules change, the physical cues maintain routine. To address staff turnover, produce a concise classroom routine guide and short video for substitutes showing how routines embed ELA opportunities (e.g., morning message writing, group storytelling). Concrete examples tying values to practice and facility features: - In a bright corner near windows (facility feature), create a multilingual book nook (value: language-rich inclusivity) where families donate books labeled in home languages and children practice retelling stories during circle time. - Use open floor space and low shelves to set up thematic play centers (value: child-centered playfulness) where children role-play bookstore or puppet theater, integrating vocabulary and print concepts; teachers rotate in to document language use for portfolios aligning with system reporting. - Install pictorial schedules at child height near classroom entry (value: safe routine consistency) so children independently move between shared-reading, choice centers, and dismissal, reducing disruptions and increasing time-on-task for ELA activities. In summary, prioritizing child-centered playfulness, language-rich inclusivity, and safe routine consistency balances the user's educational values with the realities of the existing system and facility. Anticipating conflicts---such as standardized pressures, English-only expectations, and tension between routine and flexibility---and using concrete, classroom-level strategies (integrated playful instruction, additive bilingual resources, visual routines tied to the physical space) will help resolve value conflicts while promoting affective engagement and emergent literacy outcomes. \\
  \midrule
  \textbf{Kimi-K2.6} & score 4/10 \\
  & [3-4] The report omits major issues by being incomplete, cutting off mid-sentence in Priority 2 and entirely missing Priority 3. This results in low completeness, fitting the 3-4 band, though the first section shows some logical structure. \\
  \textbf{Innospark-1T} & score 9/10 \\
  & Let me think through this carefully. \newline The question/task: You are preparing an environmental analysis report for a kindergarten English Language Arts program. Based on your analysis of the user's educational values, compatibility with existing educational systems, and the physical and social characteristics of the facility, identify and prioritize three core educational values that should guide program decisions. For each value: (a) explain why it is pr \newline My reasoning: Environmental Analysis Report: Kindergarten ELA Program \newline Methodology \newline This analysis employs an ecological systems approach, examining the intersection of educational values, institutional compatibility, and physical/social environmental factors. The prioritization framework considers: (1) developmental appropriateness for kindergarten learners, (2) alignment with ELA standards and practices, (3) facility constraints and affordances, and (4) stakeholder value systems. Three Core Educational Values \newline Value 1: Play-Based Learning \newline \textit{[\ldots\ truncated \ldots]} \newline Conclusion \newline These three values---play-based learning, inclusive multilingualism, and developmental appropriateness---form a coherent framework for kindergarten ELA program decisions. Each addresses distinct contextual factors while remaining compatible with broader educational systems when thoughtfully implemented. \\
  \specialrule{1.2pt}{8pt}{8pt}
  \multicolumn{2}{@{}l}{\textbf{38303} \quad $\Delta = +50\%$} \\
  \textbf{Question} & You are tasked with designing a multimodal learning module for an upper-level engineering course that teaches students how to create effective instructional videos, narrated demonstrations, and technical diagrams for complex system design (e.g., control systems or embedded hardware). Propose a complete project plan that a group of 3--4 students will execute over six weeks. Your plan should: (1) define clear learning objectives tied to multimodal production skills and engineering content; (2) specify the deliverables (video(s), voice-over recordings, images/diagrams, and any code or simulation artifacts) with minimum technical requirements for each; (3) outline a week-by-week schedule with milestones and checkpoints for formative feedback; (4) describe roles and responsibilities within the student team and how peer review and version control will be used; (5) detail the quality criteria and measurable success indicators you will use to decide whether the materials are effective for learner understanding (include at least two methods of formative evaluation with how and when they will be applied); and (6) produce a short storyboard and script excerpt (approx. 200--300 words) for one 6--8 minute instructional video segment that teaches a specific subsystem within the course topic. Justify your choices of multimedia elements, production techniques, and evaluation methods in terms of cognitive load, modality appropriateness, and anticipated learning gains. \\
  \textbf{Reference} & Model Answer: \newline Project Title: Multimodal Instructional Module --- Design and Demonstration of a PID-controlled DC Motor Subsystem \newline Overview and Rationale: This six-week group project teaches students to produce professional instructional videos, narrated lab demonstrations, and clear technical diagrams that communicate the design, simulation, and implementation of a PID-controlled DC motor used in an embedded control system. Objectives integrate multimodal production skills with engineering learning outcomes so students learn both content and how to teach it effectively. \newline 1) Learning Objectives \newline - Technical content objectives: \newline a. Design and simulate a PID controller for a DC motor to meet specified step-response criteria (settling time \textless{} 0.5 s, overshoot \textless{} 10\%). \newline \textit{[\ldots\ truncated \ldots]} \newline Finally, derivative action, Kd, predicts future error by measuring slope; it damps overshoot. The purple trace shows the same Kp and Ki as the red trace but with added Kd --- the peak is reduced and the response smooths out. In practice, we combine these effects: start with Kp to get near the target response, add Kd to curb overshoot, then introduce Ki carefully to remove residual steady-state error. \newline We'll now run three parameter sets in simulation and measure the resulting settling time and overshoot, then replicate the best set on hardware. Follow the on-screen checklist for safe wiring and data logging so you can reproduce these results in your lab." \\
  \midrule
  \textbf{Kimi-K2.6} & score 3/10 \\
  & [3-4] The answer is incomplete, cutting off mid-sentence in Week 2. It entirely misses sections 4, 5, and 6 (roles, quality criteria/evaluation, storyboard/script). Consequently, it fails to address accessibility, interactivity, feedback, and detailed educational design, resulting in limited learning effectiveness and low technical completeness. \\
  \textbf{Innospark-1T} & score 8/10 \\
  & Let me think through this carefully. \newline The question/task: You are tasked with designing a multimodal learning module for an upper-level engineering course that teaches students how to create effective instructional videos, narrated demonstrations, and technical diagrams for complex system design (e.g., control systems or embedded hardware). Propose a complete project plan that a group of 3--4 stud \newline My reasoning: Multimodal Learning Module Project Plan: "Embedded Control Systems Communication" \newline 1. Learning Objectives \newline Multimodal Production Skills Objectives: \newline - Create professional instructional videos using proper lighting, framing, and editing techniques \newline \textit{[\ldots\ truncated \ldots]} \newline | Storyboard Document | Scene-by-scene breakdown with timing, visual descriptions, and narration text | \newline 3. \\
  \specialrule{1.2pt}{8pt}{8pt}
  \bottomrule
\end{longtable}